\documentclass[11pt]{article}
\usepackage[review]{acl}
\usepackage{times}
\usepackage{latexsym}
\usepackage[T1]{fontenc}
\usepackage[utf8]{inputenc}
\usepackage{microtype}
\usepackage{amsmath,amssymb,amsthm}

\newtheorem{proposition}{Proposition}
\newtheorem{theorem}{Theorem}
\newtheorem{corollary}{Corollary}
\newtheorem{assumption}{Assumption}
\newcommand{\zC}{z_C}
\newcommand{\zS}{z_S}
\newcommand{\dC}{d_C}
\newcommand{\dS}{d_S}
\newcommand{\dZ}{d_Z}
\newcommand{\simf}{\operatorname{sim}}
\newcommand{\loss}{\mathcal{L}}

\title{Invariant Concept Extraction from Controlled Contrasts}
\author{}
\begin{document}
\maketitle

\begin{abstract}
\textbf{Placeholder.} We study controlled contrastive learning for separating invariant and varying information in frozen representations. The abstract will be completed once the empirical setting and results are finalized.
\end{abstract}

\section{Introduction}
\textbf{Placeholder.} Motivation, contributions, and empirical context will be added here.

\section{Related Work}
\textbf{Placeholder.} Relevant work will be added here.

\section{Method}
\textbf{Placeholder.} The implementation and data-construction details will be added here.

\section{Theoretical Analysis}
\label{sec:theory}
\subsection{Generative setting and controlled pairs}

Assume that an observation is generated by two latent blocks,
\begin{equation}
X=g(C,S), \qquad H=F(X),
\end{equation}
where $C$ denotes information intended to remain invariant under a chosen transformation relation, $S$ denotes information allowed to vary, and $F$ is a frozen representation model. We learn two maps,
\begin{equation}
Z_C=f_C(H), \qquad Z_S=f_S(H).
\end{equation}
We do not assume $C\perp S$. The objective is relative specialization: $Z_C$ should respond more strongly to changes in $C$ than to changes in $S$, and conversely for $Z_S$.

For an anchor $x_a=g(c_a,s_a)$, construct a $C$-positive $x_p=g(c_p,s_p)$ and a $C$-negative $x_n=g(c_n,s_n)$. Rather than assuming perfect interventions, we assume bounded fidelity:
\begin{align}
\dC(c_a,c_p)&\leq\eta_C, & \dS(s_a,s_p)&\geq\kappa_S,\\
\dS(s_a,s_n)&\leq\eta_S, & \dC(c_a,c_n)&\geq\kappa_C.
\end{align}
The ideal controlled setting has $\eta_C=\eta_S=0$.

\subsection{Matched contrastive objective}
For the invariant branch, let
\begin{align}
u^+&=\simf(f_C(h_a),f_C(h_p)),\\
u^-&=\simf(f_C(h_a),f_C(h_n)).
\end{align}
With temperature $\tau>0$, the matched objective is
\begin{align}
\loss_C&=-\log\frac{\exp(u^+/\tau)}{\exp(u^+/\tau)+\exp(u^-/\tau)}\\
&=\log\left(1+\exp\left(\frac{u^- - u^+}{\tau}\right)\right).
\label{eq:matched-loss}
\end{align}
For $Z_S$, we exchange the controlled contrasts: approximately equal $S$ and different $C$ is positive, while approximately equal $C$ and different $S$ is negative. The partition objective is $\loss_{\mathrm{part}}=\loss_C+\loss_S$. It requires neither reconstruction nor exact statistical independence.

Unrestricted in-batch InfoNCE is unsuitable for this purpose: it pushes down the similarity of every non-positive example, even though $c_i\ne c_j$ does not imply that two observations share no invariant information. They may share an entity, relation, topic, or another component of $C$. A matched negative instead controls $S$, so the contrast cannot be solved merely through a difference in $S$.

\begin{proposition}[Relative margin]\label{prop:margin}
If $\loss_C\leq\epsilon$ for a training triplet, then
\begin{equation}
u^+-u^-\geq\tau\log\frac{1}{\exp(\epsilon)-1}.
\end{equation}
In particular, when $\epsilon<\log 2$, $\simf(Z_C^a,Z_C^p)>\simf(Z_C^a,Z_C^n)$.
\end{proposition}
\begin{proof}
Equation~\eqref{eq:matched-loss} gives $\exp((u^--u^+)/\tau)\leq\exp(\epsilon)-1$. Taking logarithms and rearranging proves the claim.
\end{proof}

\subsection{Regularity assumptions}
Write $\zC(c,s)$ for the invariant-branch representation. In addition to bounded fidelity, assume the following.
\begin{assumption}[Observed alignment and separation]\label{ass:observed}
For observed positive pairs, $\dZ(\zC(c_a,s_a),\zC(c_p,s_p))\leq\delta_{\mathrm{obs}}$. For observed matched negatives, $\dZ(\zC(c_a,s_a),\zC(c_n,s_n))\geq\gamma_{\mathrm{obs}}$.
\end{assumption}
\begin{assumption}[Local regularity]\label{ass:lipschitz}
The representation is locally Lipschitz in each block:
\begin{align}
\dZ(\zC(c,s),\zC(c',s))&\leq L_C\dC(c,c'),\\
\dZ(\zC(c,s),\zC(c,s'))&\leq L_S\dS(s,s').
\end{align}
\end{assumption}
\begin{assumption}[Nuisance coverage]\label{ass:coverage}
For every unseen nuisance realization $s^\star$, there is a covered training realization $\tilde{s}$ with $\dS(s^\star,\tilde{s})\leq\varepsilon$. The observed alignment and separation conditions hold for the corresponding covered controlled pairs.
\end{assumption}

Pair imperfection yields the effective bounds
\begin{equation}
\delta=\delta_{\mathrm{obs}}+L_C\eta_C,\qquad \gamma=\gamma_{\mathrm{obs}}-L_S\eta_S.
\label{eq:effective-bounds}
\end{equation}

\begin{theorem}[Unseen-$S$ stability]\label{thm:stability}
Let $s^\star$ and $s^{\star\star}$ be unseen nuisance realizations with covered training realizations within $\varepsilon$. For fixed $c$,
\begin{equation}
\dZ(\zC(c,s^\star),\zC(c,s^{\star\star}))\leq\delta_{\mathrm{obs}}+L_C\eta_C+2L_S\varepsilon.
\end{equation}
\end{theorem}
\begin{proof}
Apply the triangle inequality through the two covered nuisance realizations, then Assumption~\ref{ass:lipschitz} to each unseen-to-covered segment and Equation~\eqref{eq:effective-bounds} to the middle segment.
\end{proof}

\begin{theorem}[Unseen-$C$ separation]\label{thm:separation}
For $c_i\ne c_j$ under an unseen nuisance $s^\star$,
\begin{equation}
\dZ(\zC(c_i,s^\star),\zC(c_j,s^\star))\geq\gamma_{\mathrm{obs}}-L_S\eta_S-2L_S\varepsilon.
\end{equation}
\end{theorem}
\begin{proof}
Use the reverse triangle inequality through a covered nuisance realization, then Assumption~\ref{ass:lipschitz} and Equation~\eqref{eq:effective-bounds}.
\end{proof}

\begin{corollary}[Sufficient operating condition]\label{cor:condition}
Same-$C$ observations remain closer than different-$C$ observations whenever
\begin{equation}
\gamma_{\mathrm{obs}}>\delta_{\mathrm{obs}}+L_C\eta_C+L_S\eta_S+4L_S\varepsilon.
\end{equation}
\end{corollary}
This condition is sufficient, not necessary.

\subsection{Dependence imposes unavoidable leakage}
\begin{theorem}[Leakage lower bound]\label{thm:leakage}
For discrete $C$ and any learned representation $Z_C$,
\begin{equation}
I(Z_C;S)\geq I(C;S)-H(C\mid Z_C).
\end{equation}
\end{theorem}
\begin{proof}
The chain rule gives $I(C,Z_C;S)=I(C;S)+I(Z_C;S\mid C)=I(Z_C;S)+I(C;S\mid Z_C)$. Non-negativity of $I(Z_C;S\mid C)$ and $I(C;S\mid Z_C)\leq H(C\mid Z_C)$ establish the result.
\end{proof}

If $C$ is decoded from $Z_C$ with error probability $p_e$, Fano's inequality yields
\begin{equation}
I(Z_C;S)\geq I(C;S)-h(p_e)-p_e\log(|\mathcal{C}|-1).
\end{equation}
Thus, when $C$ and $S$ are dependent, faithful preservation of $C$ entails a nonzero leakage floor. The appropriate goal is to remove $S$ information beyond that intrinsically coupled to $C$, rather than to require $I(Z_C;S)=0$.

\subsection{Propagation to sparse autoencoders}
Let $A=\phi(Z_C)$ be an SAE encoding. For a Top-$k$ SAE, suppose locally that $\|\phi(z)-\phi(z')\|\leq L_{\mathrm{SAE}}\|z-z'\|$ while its selected support is fixed. Theorem~\ref{thm:stability} then gives
\begin{equation}
\begin{aligned}
&\|\phi(\zC(c,s^\star))\\
&\quad{}-\phi(\zC(c,s^{\star\star}))\|\\
&\leq L_{\mathrm{SAE}}\bigl(\delta_{\mathrm{obs}}+L_C\eta_C+2L_S\varepsilon\bigr).
\end{aligned}
\end{equation}
Moreover, if ordered SAE preactivations have Top-$k$ support margin $\rho=q_{(k)}-q_{(k+1)}$, the support is unchanged whenever $\|W\Delta\zC\|_\infty<\rho/2$. A sufficient condition is
\begin{equation}
\begin{aligned}
\|W\|_{\infty\leftarrow 2}(\delta_{\mathrm{obs}}+L_C\eta_C
&+2L_S\varepsilon)<\rho/2.
\end{aligned}
\end{equation}

\paragraph{Scope.} The analysis establishes controlled relative preference, bounded nuisance drift under the stated assumptions, and downstream feature stability under a support-margin condition. It does not prove $Z_C=C$, $Z_S=S$, $Z_C\perp S$, globally identifiable factors, human-interpretable concepts, or generalization to arbitrarily distant nuisance variation.

\section{Experiments}
\textbf{Placeholder.} Experimental design and results will be added here.

\section{Conclusion}
\textbf{Placeholder.} The conclusion will be added here.
\end{document}
